\originalsection{逐讲内容与原件可得性}
MAS二十讲均有实际PDF. 下表列出正文推导而非仅出现的名词; 完整研究定理的证明层级仍按前述边界.
\begin{longtable}{p{.09\linewidth}p{.12\linewidth}p{.69\linewidth}}
\toprule 讲 & 章 & 实际写入的推导/协议\\\midrule\endhead
1 & 5 & Kraus、环境等距、Choi与完全正；根保真度\\
2 & 6 & 连续误差展开、Knill--Laflamme必要与充分结构、Shor码\\
3 & 6 & CSS对偶码、Fourier支持、Steane综合征\\
4 & 7 & 二元辛表示、投影维数、逻辑正规化子、测量更新\\
5 & 7 & 环面边/星算符、同调逻辑环、距离与anyon相位\\
6 & 7 & Clifford更新、横向传播与编码逻辑操作；PS2中途错误\\
7 & 8 & 验证猫、四轮停机、Y资源注入与级联递推\\
8 & 9 & 图态稳定子、删除Z测量、链式传态与前馈\\
9 & 4 & Bell传态分支与UP副产物；第9章一步测量门\\
10 & 9 & 绝热搜索隙与调度、历史态传播项、合法钟和能量保护\\
11 & 3 & 连续搜索、行走色散、Ambainis缓存与五维矩阵、成本优化\\
12 & 10 & Abel陪集振幅、离散对数显式函数、Pell调节子周期流程\\
13 & 10 & 表示正交与矩阵QFT、二面体相位合成、Heisenberg实例\\
14 & 11 & BSC容量、强/弱典型、CPTP压缩器、保真度界、类型集中\\
15 & 12 & Gram平方根POVM、相对熵条件、Holevo、HSW打包路线\\
16 & 12 & 正规化容量优化、擦除信道、参考系统与父/母协议账\\
17 & 12 & 主序见第11章；混合Bell谱/PPT阈值；非局域CNOT与SWAP\\
18 & 13 & 多项式指纹、SMP模型、量子指纹误判与一次性签名全流程\\
19 & 13 & 策略/Nash、PQ、囚徒收益表、SU2差别、公共品GHZ机制\\
20 & 13 & EPR抽样、纯化保真度到熵、CSS到BB84、移除量子存储\\
\bottomrule\end{longtable}
18.435实际取得的原讲PDF仅L2、4、5、11、12、13、14、16、19、23十份. L2的基础量子数学、L19的详细实验器件作为先修或与AMO卷衔接, 其余八讲的计算/纠错链见上列章节. L1、3、6--10、15、17--18、20--22没有取得原讲文件, 对应主题可由MAS和作业材料建立, 但不称那些缺失原讲已经翻译. 完整逐讲文件链接及状态见\texttt{lecture-coverage.json}.
